\documentclass{article}
\usepackage{amsmath, amsthm}

\newtheorem{theorem}{Theorem}
\newtheorem{definition}{Definition}

\title{A Sample Paper with SemWeave Annotations}
\author{Example Author}

\begin{document}
\maketitle

% mcp: begin region role=section name=introduction anchors=[sec:intro]
\section{Introduction}

This paper demonstrates how SemWeave annotations can be used to provide
structured navigation over a LaTeX document. The key idea is that structure
is declared through comments, not inferred from LaTeX commands.

We will define important concepts in Section~\ref{sec:background} and
present our main result in Section~\ref{sec:main}.
% mcp: end

% mcp: begin region role=section name=background anchors=[sec:background]
\section{Background}

% mcp: begin region role=definition name=graph-def anchors=[def:graph]
\begin{definition}[Graph]
A graph $G = (V, E)$ consists of a set of vertices $V$ and a set of
edges $E \subseteq V \times V$.
\end{definition}
% mcp: end

% mcp: begin region role=definition name=tree-def anchors=[def:tree]
\begin{definition}[Tree]
A tree is a connected acyclic graph. Every tree with $n$ vertices has
exactly $n-1$ edges. See also Definition~\ref{def:graph}.
\end{definition}
% mcp: end

% mcp: end

% mcp: begin region role=section name=main-results anchors=[sec:main]
\section{Main Results}

% mcp: begin region role=theorem name=main-theorem anchors=[thm:main]
\begin{theorem}[Main Theorem]
Every connected graph contains a spanning tree.
\end{theorem}
% mcp: end

% mcp: begin region role=proof name=main-proof anchors=[prf:main]
\begin{proof}
We proceed by induction on the number of edges. The base case is
immediate from Definition~\ref{def:tree}. For the inductive step,
if $G$ contains a cycle, remove any edge from the cycle. The resulting
graph is still connected and has fewer edges. By induction, it contains
a spanning tree, which is also a spanning tree of $G$.
\end{proof}
% mcp: end

% mcp: begin region role=example name=example1 anchors=[ex:triangle]
\subsection{Example}
Consider the complete graph $K_3$ (triangle). It has 3 vertices and 3 edges.
Removing any single edge yields a spanning tree (a path of length 2).
This illustrates Theorem~\ref{thm:main}.
% mcp: end

% mcp: end

\end{document}
